\chapterpage{sources}{两条拟合与实验输入的出处}
\textbf{Celliers 铝。} 速度采用 km/s，初始密度取材料库中的 $\rho_0=2.70$ g/cm³。拟合关系为
\[
 U_s=17.992+1.167(u_p-9.838).
\]
文件在 0 到 50 km/s 的粒子速度范围内采样 500 点，拟合系数及其不确定度见表头。\href{https://doi.org/10.1063/1.2140077}{Celliers 等，J. Appl. Phys. 98, 113529（2005）。}

\textbf{Hamel α-石英。} 初始密度取材料库中的 $\rho_0=2.65$ g/cm³。低于 15 km/s 时使用二次式，自连接点起沿切线延拓。文件在 0.1 到 50 km/s 范围内采样 500 点。
\[
\begin{array}{ll}
U_s=2.40+1.68u_p-0.0155u_p^2 & (u_p<15),\\
U_s=5.8875+1.215u_p & (u_p\geq15).
\end{array}
\]
这条石英标准曲线见 Hamel 等关于 CH$_{1.36}$ 的论文。文件的采样区间表示软件提供的数据范围，其他工况下的适用性应对照原文的校准数据。\href{https://doi.org/10.1103/PhysRevB.86.094113}{Hamel 等，Phys. Rev. B 86, 094113（2012）。}

\subhead{第 29400 号实验的两个输入}
Hicks 等人表 I 给出的实测冲击波速度为铝 $18.11\pm0.32$、石英 $17.96\pm0.26$ µm/ns。1 µm/ns 等于 1 km/s，输入时数值不变。论文利用其采用的铝标准推得石英压强 4.79 Mbar、密度 6.03 g/cm³、粒子速度 10.07 km/s。本练习取两个实测波速作为输入。\href{https://doi.org/10.1063/1.2009528}{Hicks 等，Phys. Plasmas 12, 082702（2005），表 I。}

\subhead{独立核算单位}
在 $u_p=10.0$ km/s 处，Hamel 石英表给出 $U_s=17.65$ km/s，压强为 $0.01\times2.65\times17.65\times10=4.67725$ Mbar。相同粒子速度代入铝的解析式得到 4.90888 Mbar，软件对铝表采样点线性插值得到 4.90894 Mbar。
